\documentclass[11pt]{article}
\usepackage[utf8]{inputenc}
\usepackage[spanish]{babel}
\usepackage{amsmath,amssymb,amsthm}
\usepackage{graphicx}
\usepackage[margin=2.5cm]{geometry}
\usepackage{enumitem}

\newtheorem*{definicion}{Definición}
\newtheorem*{ejemplo}{Ejemplo}
\newtheorem*{teorema}{Teorema}
\newtheorem*{observacion}{Observación}

\title{Semana 8 -- Materia}
\author{Tutorías -- Introducción al Cálculo}
\date{28/9--2/10/2026}

\begin{document}
\maketitle

\section{Funciones inyectivas}

Hasta ahora sabemos que una función asigna a cada elemento de su dominio una
única imagen. Nos interesa ahora el caso en que, además, elementos distintos
del dominio tienen imágenes distintas: esto es lo que permite, más adelante,
``deshacer'' la función y definir una inversa.

\begin{definicion}[Función inyectiva]
Sea $f$ función. $f$ es inyectiva o uno a uno si y solo si
\[
\forall x \in \mathsf{Dom}\,f\ \forall y \in \mathsf{Dom}\,f\ (x \neq y \to f(x) \neq f(y)).
\]
\end{definicion}

Esta definición es equivalente a su forma contrarrecíproca, que en la práctica
es la que más se usa para demostrar inyectividad:
\[
\forall x \in \mathsf{Dom}\,f\ \forall y \in \mathsf{Dom}\,f\ (f(x) = f(y) \to x = y).
\]

Gráficamente, $f$ es inyectiva cuando toda recta horizontal intersecta su
gráfico en a lo más un punto.

\begin{ejemplo}
Sea $f(x) = \dfrac{x+1}{x-2}$, con $\mathsf{Dom}\,f = \mathbb{R}-\{2\}$.
Sean $x,y \in \mathbb{R}-\{2\}$ con $f(x)=f(y)$. Entonces
\[
\frac{x+1}{x-2} = \frac{y+1}{y-2} \iff (x+1)(y-2) = (y+1)(x-2) \iff xy-2x+y-2 = xy-2y+x-2,
\]
de donde $3y = 3x$, es decir $x=y$. Por lo tanto $f$ es inyectiva.
\end{ejemplo}

\begin{ejemplo}
Sea $f(x) = 2x+1$. Si $f(x) = f(y)$ entonces $2x+1 = 2y+1$, de donde $x=y$.
Luego $f$ es inyectiva. En cambio, $g(x) = x^2$ no es inyectiva, pues
$g(1) = g(-1) = 1$ con $1 \neq -1$.
\end{ejemplo}

\section{Relación inversa}

Toda función, sea o no inyectiva, tiene asociada una relación inversa (que no
siempre es función).

\begin{definicion}[Relación inversa]
Sea $f$ función real. La relación inversa $f^{-1}$ de $f$ se define por
\[
f^{-1} = \{(y,x) \in \mathsf{Rec}\,f \times \mathsf{Dom}\,f : f(x) = y\}.
\]
\end{definicion}

Intuitivamente, si elementos distintos del dominio tienen imágenes distintas
(es decir, si $f$ es inyectiva), entonces cada elemento del recorrido tiene
una única preimagen, y por lo tanto $f^{-1}$ resulta ser función. Este es el
contenido del siguiente teorema.

\begin{teorema}
Sea $f$ función. Entonces $f$ es inyectiva si y solo si $f^{-1}$ es función.
\end{teorema}

\emph{Idea de la demostración.} Si $f$ es inyectiva y $(y,x), (y,x') \in f^{-1}$,
entonces $f(x) = y = f(x')$, y como $f$ es inyectiva, $x=x'$; luego $f^{-1}$ es
función. Recíprocamente, si $f^{-1}$ es función y $f(x) = f(x') = y$, entonces
$(y,x), (y,x') \in f^{-1}$, y como $f^{-1}$ es función, $x=x'$; luego $f$ es
inyectiva.

\begin{observacion}
Cuando $f^{-1}$ es función, se cumple
\[
\mathsf{Dom}\,f^{-1} = \mathsf{Rec}\,f, \qquad \mathsf{Rec}\,f^{-1} = \mathsf{Dom}\,f,
\]
y además, si $x \in \mathsf{Dom}\,f$ e $y \in \mathsf{Rec}\,f$, entonces
\[
f(x) = y \leftrightarrow f^{-1}(y) = x.
\]
\end{observacion}

\begin{ejemplo}
Retomando $f(x) = \dfrac{x+1}{x-2}$, ya sabemos que $f$ es inyectiva, con
$\mathsf{Dom}\,f = \mathbb{R}-\{2\}$ y $\mathsf{Rec}\,f = \mathbb{R}-\{1\}$.
Para $y \in \mathbb{R}-\{1\}$, buscamos $x$ tal que $f(x) = y$:
\[
\frac{x+1}{x-2} = y \iff x+1 = y(x-2) \iff x(1-y) = -1-2y \iff x = \frac{2y+1}{y-1}.
\]
Luego $f^{-1}$ es la función real con dominio $\mathbb{R}-\{1\}$, recorrido
$\mathbb{R}-\{2\}$ y definida por $f^{-1}(y) = \dfrac{2y+1}{y-1}$.
\end{ejemplo}

\section{Biyecciones}

Decimos que una función de $A$ en $B$ es \emph{sobre} $B$ cuando su recorrido
es exactamente $B$. Combinando esto con la inyectividad obtenemos el concepto
de biyección.

\begin{definicion}[Biyección de $A$ sobre $B$]
Sean $A,B$ conjuntos. $f$ es biyección de $A$ sobre $B$ si y solo si $f$ es
función inyectiva de $A$ en $B$ y $\mathsf{Rec}\,f = B$.
\end{definicion}

\begin{ejemplo}
Sea $f:\mathbb{R}\to\mathbb{R}$ definida por $f(x) = 2x$.
\begin{itemize}
\item $f$ es uno a uno: si $f(x)=f(y)$ entonces $2x=2y$, luego $x=y$.
\item $f$ es sobre $\mathbb{R}$: dado $y \in \mathbb{R}$, sea $x = y/2 \in \mathbb{R}$;
entonces $f(x) = 2(y/2) = y$, luego $y \in \mathsf{Rec}\,f$.
\end{itemize}
Por lo tanto, $f$ es biyección de $\mathbb{R}$ sobre $\mathbb{R}$.
\end{ejemplo}

\begin{ejemplo}
Sea $f:\mathbb{N}\to\mathbb{N}$ definida por $f(x) = 3x$.
\begin{itemize}
\item $f$ es uno a uno: si $f(x)=f(y)$ entonces $3x=3y$, luego $x=y$.
\item $f$ no es sobre $\mathbb{N}$: $2 \notin \mathsf{Rec}\,f$, pues si $2=f(x)$ para algún
$x\in\mathbb{N}$, entonces $x = 2/3 \notin \mathbb{N}$.
\end{itemize}
Por lo tanto, $f$ no es biyección de $\mathbb{N}$ sobre $\mathbb{N}$.
\end{ejemplo}

\section{Propiedades de las biyecciones}

\begin{teorema}
Si $f$ es función inyectiva, entonces $f^{-1}$ también es función inyectiva.
\end{teorema}

\emph{Idea de la demostración.} Como $f$ es inyectiva, $f^{-1}$ es función
(por el teorema anterior). Además,
\[
f^{-1}(x) = f^{-1}(y) \implies f(f^{-1}(x)) = f(f^{-1}(y)) \implies x=y,
\]
así que $f^{-1}$ también es inyectiva.

\begin{teorema}
Si $f$ es biyección de $A$ sobre $B$, entonces $f^{-1}$ es biyección de $B$
sobre $A$.
\end{teorema}

\emph{Idea de la demostración.} Por el teorema anterior, $f^{-1}$ es función
inyectiva, y $\mathsf{Dom}\,f^{-1} = \mathsf{Rec}\,f = B$. Falta ver que
$f^{-1}$ es sobre $A$: dado $y \in A$, como $f$ es sobre $B$, existe $x \in B$
con $f^{-1}(x) = y$ (usando que $f(y)=x$), de donde $y \in \mathsf{Rec}\,f^{-1}$.

\begin{teorema}
Si $f$ es biyección de $A$ sobre $B$ y $g$ es biyección de $B$ sobre $C$,
entonces $g\circ f$ es biyección de $A$ sobre $C$.
\end{teorema}

\emph{Idea de la demostración.} Se verifica que $\mathsf{Dom}(g\circ f) = A$ y
que $\mathsf{Rec}(g\circ f) = C$ (usando que $f$ y $g$ son sobre). Para la
inyectividad, si $g\circ f(x) = g\circ f(x')$, entonces $g(f(x)) = g(f(x'))$;
como $g$ es uno a uno, $f(x)=f(x')$, y como $f$ es uno a uno, $x=x'$.

\section{Aplicación: funciones por tramos}

Cuando una función está definida por tramos, probar inyectividad y
sobreyectividad requiere analizar cada tramo y también las combinaciones entre
tramos distintos, como muestra el siguiente problema.

\begin{ejemplo}[Problema 3.1 del libro]
Sean $f,g$ funciones reales definidas por
\[
f(x) = \begin{cases} x^2-1 & \text{si } x>0 \\ 2x-1 & \text{si } x \le 0 \end{cases}
\qquad
g(x) = \begin{cases} x^2-1 & \text{si } x>-1 \\ 2x-1 & \text{si } x \le -1. \end{cases}
\]
Se pide mostrar que $f$ es biyección sobre $\mathbb{R}$, que $g$ no lo es, y
determinar $g \circ f^{-1}$.

\emph{Inyectividad de $f$.} Como $f$ está definida por tramos, hay que
considerar tres casos para $x \neq y$:
\begin{itemize}
\item $x,y>0$: $x^2 \neq y^2$, luego $f(x) = x^2-1 \neq y^2-1 = f(y)$.
\item $x,y\le 0$: $2x \neq 2y$, luego $f(x) = 2x-1 \neq 2y-1 = f(y)$.
\item $x\le 0 < y$: $f(x) = 2x-1 \le -1$ y $f(y) = y^2-1 > -1$, luego $f(x)\neq f(y)$.
\end{itemize}
En los tres casos $f(x)\neq f(y)$, así que $f$ es inyectiva.

\emph{Sobreyectividad de $f$.} Dado $y\in\mathbb{R}$, distinguimos $y\le -1$ e
$y>-1$:
\begin{itemize}
\item Si $y \le -1$, sea $x = \frac{y+1}{2} \le 0$; entonces $f(x) = 2\left(\frac{y+1}{2}\right)-1 = y$.
\item Si $y > -1$, sea $x = \sqrt{y+1} > 0$; entonces $f(x) = (\sqrt{y+1})^2-1 = y$.
\end{itemize}
Luego $f$ es sobre $\mathbb{R}$, y por lo tanto $f$ es biyección de $\mathbb{R}$
sobre $\mathbb{R}$.

\emph{$g$ no es biyección.} $g$ no es inyectiva: con $x_1=-\frac12$, $x_2=\frac12$
(ambos en el tramo $x>-1$), $g(x_1) = \frac14-1 = -\frac34 = g(x_2)$ pese a que
$x_1\neq x_2$. Tampoco es sobre $\mathbb{R}$ (por ejemplo, no existe $x$ con
$g(x)=-2$).

\emph{Cálculo de $f^{-1}$.} Aprovechando el trabajo de la sobreyectividad,
\[
f^{-1}(y) = \begin{cases} \dfrac{y+1}{2} & \text{si } y \le -1 \\[4pt] \sqrt{y+1} & \text{si } y > -1. \end{cases}
\]

\emph{Cálculo de $g\circ f^{-1}$.} Distinguiendo según el valor de $x$ (y
usando que $f^{-1}(x) > -1 \leftrightarrow x > -3$ cuando $x \le -1$), se
obtiene:
\[
g\circ f^{-1}(x) = \begin{cases} x & \text{si } x \le -3 \lor x > -1 \\[4pt] \left(\dfrac{x+1}{2}\right)^2-1 & \text{si } -3 < x \le -1. \end{cases}
\]
\end{ejemplo}

\end{document}
